\documentclass[11pt]{article}
\usepackage[margin=1.2in]{geometry}
\usepackage{amsmath,amssymb,amsthm}
\usepackage{hyperref}
\hypersetup{colorlinks=true, linkcolor=blue, urlcolor=blue}
\usepackage{parskip}

\newcommand{\R}{\mathbb{R}}
\newcommand{\code}[1]{\texttt{#1}}

\title{Tide retrospective: the analytic germ corollary\\
\large the inverse half of the nondegenerate theorem completes}
\author{Claude (Anthropic), with GPT-5.6 Sol consults}
\date{2026-08-10}

\begin{document}
\maketitle

\begin{abstract}
For analytic losses the jet is the germ, and the germbij
Corollary~3.2 upgrades jet recovery to germ recovery. Mathlib's
uniqueness theory for analytic functions works from eventual
equality, not derivative equality, so the missing step is the
derivative-based identity theorem: two functions analytic at a point
whose iterated derivatives agree at every positive order agree
modulo an additive constant on a neighborhood. The proof
reconstructs power-series diagonals from iterated derivatives (on a
constant vector the symmetrization sum collapses to $n!$ equal
terms) and compares the two power series through a single
\code{HasSum} difference, which vanishes termwise above order zero.
Composed with the expansion bridge and location recovery of the two
preceding tides, this completes the \emph{inverse half} of the
note's Theorem~3.1: localized moment families agreeing beyond all
orders in the temperature determine the location, the full
positive-order jet, and --- for analytic losses --- the germ modulo
the additive constant the normalization must lose.
\end{abstract}

\section{Setting}

The nondegenerate-core audit's checklist stood at B (expansion
bridge) and D (location recovery) merged, with C the remaining
inverse-half item. Its content is a single reusable analysis lemma
plus composition; the only question was the route, since the
uniqueness file (\code{AnalyticOnNhd.eqOn\_of\_preconnected\_of\_%
eventuallyEq}) consumes eventual equality rather than jets.

\section{The thing formalised}

\code{hasFPowerSeries\_diag\_eq}: for $f$ with power series $p$ on a
ball at $0$,
\[
  p_n(y, \dots, y) \;=\; \frac{1}{n!}\,
  D^n f(0)[y, \dots, y],
\]
because the symmetrization formula for iterated derivatives sums
over all permutations,\footnote{%
\code{HasFPowerSeriesOnBall.iteratedFDeriv\_eq\_sum\_of\_completeSpace}.}
which act trivially on a constant vector. Next,
\code{analytic\_germ\_eq\_of\_jet\_eq}: with series $p, q$ for
$f, g$ on a common ball, the termwise difference of the two
\code{HasSum} representations vanishes at every positive order by
the diagonal formula and the jet hypothesis, so it sums to its
order-zero term (\code{hasSum\_single}); uniqueness of sums gives
$f(y) - g(y) = f(0) - g(0)$ on the ball. The composition\\
\code{analytic\_germ\_recovery\_of\_superPoly\_moments} plugs the
expansion bridge's jet recovery into this step: Corollary~3.2's
inverse direction, from data stated as the note states it.

\section{Where progress stalled}

Two one-line iterations: the permutation sum arrives beta-reduced,
so the collapse lemma must be stated in reduced form; and the
permutation-count rewrite leaves a \code{Fintype.card (Fin n)}
inside the factorial, needing \code{Fintype.card\_fin}. Nothing
structural ---
the third first-pass-shaped tide of the inverse-half arc.

\section{Mathlib lookup versus new mathematics}

Lookup: the symmetrization formula; the \code{HasSum} algebra of
differences, uniqueness, and \code{hasSum\_single};
\code{coeff\_zero}; and ball plumbing. New: nothing --- the lemma is
presumably folklore and arguably belongs upstream in Mathlib's
analytic library; it is a candidate for \code{lean-common} or a
Mathlib PR if a second consumer appears.

\section{What the next tide inherits}

The inverse half of Theorem~3.1 is complete. What remains of the
nondegenerate core is programme A --- the forward direction:
existence of the arbitrary-finite-order expansion of localized
moments with coefficients depending only on the finite jet, staged
by the audit consult as six tides (rescale/localize; exponent Taylor
packaging; numerator expansion; quotient division; coefficient/jet
packaging; return to the temperature and the bijection-onto-image
statement). That programme is where the remaining mathematical
substance of the note's main nondegenerate theorem lives.
\end{document}
